\section{Methodologies for Enhancing Reasoning Capabilities}
\label{sec:intro-objectives}

% WHAT TO WRITE -------------------------------------------------------
% The concrete questions you aim to answer and the objectives for this
% stage of the study.
% ---------------------------------------------------------------------

Improving reasoning involves specific techniques typically applied after or on top of the conventional pipeline.

\subsection{Inference-Time Compute Scaling}

This method improves performance at the moment of the user prompt without modifying the model’s weights. It trades increased computation for higher accuracy using:

\begin{itemize}
  \item Chain-of-Thought prompting
  \item Advanced sampling and voting procedures to identify the best candidate solution.
\end{itemize}
  
\subsection{Reinforcement Learning (RL) for Reasoning}

Unlike RLHF, which relies on subjective human judgment, RL for reasoning uses automated or environment-based reward signals.

\begin{itemize}
  \item \textbf{Verification:} Rewards are given for verifiable correctness (e.g., a correct math answer or functional code).
  \item \textbf{Weight Updates:} This process modifies the model's weights through trial and error.
\end{itemize}

\subsection{Knowledge Disitllation}

This involves transferring reasoning patterns from a large, high-performing model ("teacher") to a smaller, more efficient "student" model. This is typically achieved via SFT using high-quality reasoning datasets generated by the stronger model.

% \end{itemize}
% \begin{enumerate}
%   \item \textbf{RQ1.} First research question.
%   \item \textbf{RQ2.} Second research question.
% \end{enumerate}

% \noindent Objectives:
% \begin{itemize}
%   \item First objective.
%   \item Second objective.
% \end{itemize}
